\documentclass[11pt]{amsart}

\usepackage[T1]{fontenc}
\usepackage{lmodern}
\usepackage{microtype}
\usepackage{amsmath,amssymb,amsthm}
\usepackage[hidelinks]{hyperref}

\hypersetup{
  pdfsubject={Möbius uncertainty on locally finite posets},
  pdfkeywords={Möbius inversion, uncertainty principle, locally finite poset, incidence algebra}
}

\newtheorem{theorem}{Theorem}
\newtheorem{corollary}[theorem]{Corollary}
\theoremstyle{remark}
\newtheorem{remark}[theorem]{Remark}

\newcommand{\supp}{\operatorname{supp}}

\title[$\mathcal H_1$ and $\mathcal H_2$ do not characterize $\mathcal M$]
{$\mathcal H_1$ and $\mathcal H_2$ Do Not Characterize Möbius Uncertainty}

\date{}

\subjclass[2020]{Primary 06A07; Secondary 11A25}
\keywords{Möbius inversion, uncertainty principle, locally finite poset,
incidence algebra}

\begin{document}

\begin{abstract}
Sahay left open, while writing that he considered it ``quite unlikely'',
whether the Möbius uncertainty property $\mathcal M$ of a locally finite
poset with a minimum is characterized by the two necessary conditions
$\mathcal H_1$ and $\mathcal H_2$. We confirm his expectation: the
characterization fails. For every $r\geq 1$, we construct a countable
locally finite poset $P_r$ with a minimum such that $P_r$ satisfies $\mathcal H_k$ exactly for
$1\leq k\leq r$, but fails $\mathcal M$. Thus no finite family of the
conditions $\mathcal H_k$ implies Möbius uncertainty.
\end{abstract}

\maketitle

Let $P$ be a locally finite poset with minimum $\widehat 0$. For
$f\colon P\to\mathbb C$, define
\[
  (Zf)(x)=\sum_{y\leq x}f(y),
  \qquad
  \supp(f)=\{x\in P:f(x)\neq 0\}.
\]
The poset $P$ has the \emph{Möbius uncertainty property} $\mathcal M$ if,
for every nonzero $f$, at least one of $\supp(f)$ and $\supp(Zf)$ is
infinite. For $k\geq 1$, it has property $\mathcal H_k$ if every
$k$-element set $S\subseteq P$ contains some $z\in S$ for which there are
infinitely many $x\in P$ such that
\[
  x\geq z
  \quad\text{and}\quad
  x\ngeq y\quad (y\in S\setminus\{z\}).
\]

Pollack proved the arithmetic uncertainty principle \cite{Pollack}, and
Goh formulated and studied its poset analogue \cite{Goh}. Sahay introduced
the properties $\mathcal H_k$ and proved
\[
  \mathcal H_{\mathbb N}\Longrightarrow\mathcal M
  \Longrightarrow \mathcal H_1\text{ and }\mathcal H_2,
\]
where $\mathcal H_{\mathbb N}$ means that every $\mathcal H_k$ holds
\cite[Theorems 1.2 and 1.5]{Sahay}. After showing that
$\mathcal M\nRightarrow\mathcal H_3$ \cite[Theorem 1.11]{Sahay}, he wrote,
immediately following that theorem \cite[p.~5]{Sahay}:
\begin{quote}
``This leaves open the possibility that
\[
  P\text{ has }\mathcal M\iff P\text{ has }\mathcal H_1\text{ and }
  \mathcal H_2,
\]
for general posets, but I consider this quite unlikely.''
\end{quote}

\begin{theorem}\label{thm:main}
For every integer $r\geq 1$, there is a countable locally finite poset
$P_r$ with a minimum such that
\[
  P_r\text{ has }\mathcal H_k
  \quad\Longleftrightarrow\quad
  1\leq k\leq r,
  \qquad
  P_r\text{ does not have }\mathcal M.
\]
\end{theorem}

\begin{proof}
Let $\mathcal F(\mathbb N)$ be the poset of finite subsets of
$\mathbb N=\{1,2,\ldots\}$ ordered by inclusion. For $1\leq i\leq r$, let
\[
  Q_i=\{A_i:A\in\mathcal F(\mathbb N)\}
\]
be a copy of this poset. Define
\[
  P_r=\{\widehat 0\}\sqcup\bigsqcup_{i=1}^r Q_i
\]
by placing $\widehat 0$ below every element, retaining inclusion within each
$Q_i$, and imposing no comparabilities between distinct copies. The poset is
countable and locally finite, since
\[
  [A_i,B_i]=\{C_i:A\subseteq C\subseteq B\},
  \qquad
  [\widehat 0,B_i]=\{\widehat 0\}\cup\{C_i:C\subseteq B\}.
\]

Fix $1\leq k\leq r$ and let $S\subseteq P_r$ have size $k$. If
$\widehat 0\in S$, then the other $k-1$ elements occupy at most $r-1$
copies. Choose a copy $Q_i$ disjoint from $S\setminus\{\widehat 0\}$. Taking
$z=\widehat 0$, every element of $Q_i$ is an exclusive upper witness for
$z$.

Suppose $\widehat 0\notin S$. Choose an occupied copy $Q_i$, put
$T=S\cap Q_i$, and choose $A_i\in T$ minimal under inclusion. For every
$B_i\in T\setminus\{A_i\}$, choose $b_B\in B\setminus A$, and let
\[
  E=\{b_B:B_i\in T\setminus\{A_i\}\}.
\]
For each $c\in\mathbb N\setminus(A\cup E)$, set
$X_c=(A\cup\{c\})_i$. Then $X_c\geq A_i$, while $B_i\nleq X_c$ for every
$B_i\in T\setminus\{A_i\}$; elements of $S\setminus T$ lie in other copies
and are incomparable with $X_c$. Since there are infinitely many choices of
$c$, this proves $\mathcal H_k$.

Now let $k\geq r+1$ and extend
\[
  S_0=\{\widehat 0,\varnothing_1,\ldots,\varnothing_r\}
\]
to a $k$-element set $S\subseteq P_r$. For $z=\widehat 0$, the only element
above $z$ that is not above any other element of $S$ is $\widehat 0$ itself,
since every nonminimum element lies above some $\varnothing_i$. For
$z\neq\widehat 0$, every element above $z$ is also above $\widehat 0$.
Thus $S$ has no element with infinitely many exclusive upper witnesses, and
$P_r$ fails $\mathcal H_k$.

Finally, define $f\colon P_r\to\mathbb C$ by
\[
  f(\widehat 0)=1,
  \qquad
  f(\varnothing_i)=-1\quad(1\leq i\leq r),
  \qquad
  f(x)=0\quad\text{otherwise}.
\]
For every $A_i\in Q_i$,
\[
  (Zf)(A_i)=f(\widehat 0)+f(\varnothing_i)=0,
  \qquad
  (Zf)(\widehat 0)=1.
\]
Hence both $f$ and $Zf$ are nonzero and finitely supported, so $P_r$ does
not have $\mathcal M$.
\end{proof}

\begin{corollary}
The possibility left open in \cite[p.~5]{Sahay} is false. More generally, no
finite collection
of the conditions $\mathcal H_k$ implies the Möbius uncertainty property for
locally finite posets with a minimum.
\end{corollary}

\begin{proof}
The first assertion is Theorem~\ref{thm:main} with $r=2$. Given a finite set
$K\subseteq\mathbb N$, take $r\geq\max(K\cup\{1\})$; then $P_r$ has every
$\mathcal H_k$ with $k\in K$ but does not have $\mathcal M$.
\end{proof}

\begin{remark}
The displayed witness is support-minimal. Indeed, if
$0\neq h\colon P_r\to\mathbb C$ has $s=|\supp(h)|\leq r$, then
$\mathcal H_s$ supplies infinitely many $x$ for which
$(Zh)(x)=h(z)\neq 0$ for some $z\in\supp(h)$. Thus every finitely supported
witness to failure of $\mathcal M$ has support at least $r+1$, and the
function $f$ above attains this bound.

In $P_r$ the
copies $Q_1,\ldots,Q_r$ are the up-sets of
$\varnothing_1,\ldots,\varnothing_r$ and partition
$P_r\setminus\{\widehat 0\}$. What makes $\mathcal H_1,\ldots,\mathcal H_r$
survive is the number of copies rather than the partition alone: for
$\widehat 0\in S$ it is the pigeonhole over exactly $r$ copies that leaves a
copy free, while for $\widehat 0\notin S$ the witnesses are produced inside
a single copy, for a set of any size. Among the sets of size $r+1$, the only
one with no element having infinitely many exclusive upper witnesses is
$S_0=\{\widehat 0,\varnothing_1,\ldots,\varnothing_r\}$, and the failure of
$\mathcal H_k$ for $k\geq r+1$ comes from the extensions of $S_0$.
For $r\geq 2$, the poset $P_r$ is not
a join-semilattice, since elements in distinct copies have no common upper
bound; hence the construction lies outside Sahay's join-semilattice theorem
\cite[Theorem 1.10 and Section 4]{Sahay}.
\end{remark}

\begin{thebibliography}{9}

\bibitem{Goh}
M.~K. Goh,
\emph{An uncertainty principle for Möbius inversion on posets},
Contrib. Discrete Math. \textbf{21} (2026), no.~1, 142--148,
\href{https://doi.org/10.55016/6ej3y980}{doi:10.55016/6ej3y980}.

\bibitem{Pollack}
P. Pollack,
\emph{The Möbius transform and the infinitude of primes},
Elem. Math. \textbf{66} (2011), no.~3, 118--120,
\href{https://doi.org/10.4171/EM/178}{doi:10.4171/EM/178}.

\bibitem{Sahay}
A. Sahay,
\emph{A note on the Möbius uncertainty principle for posets},
\href{https://arxiv.org/abs/2603.01018}{arXiv:2603.01018v2}, 4 May 2026.

\end{thebibliography}

\end{document}